% ============================================================
% Inhalt: Paare, Photonen und Löcher
% Torsionserhaltung als gemeinsames Prinzip
% hinter Cooper-Paar, Exziton und SPDC
% März 2026
%
% Konvention dieses Dokuments:
%   [STANDARD] = etablierte Physik mit Quellen
%   [FFGFT]    = Interpretation im FFGFT-Rahmen
% ============================================================

\providecommand{\ket}[1]{\left|#1\right\rangle}
\providecommand{\bra}[1]{\left\langle#1\right|}

% ============================================================
\section{Ausgangsfrage und Methodik}
\label{sec:ausgangsfrage}
% ============================================================

Drei Phänomene der modernen Physik zeigen auf den ersten Blick
eine strukturelle Ähnlichkeit:

\begin{itemize}
  \item Ein Photon spaltet sich in zwei Photonen niedrigerer Frequenz auf
    (parametrische Fluoreszenz, SPDC).
  \item Ein angeregtes Elektron bildet mit einem zurückbleibenden
    Loch ein gebundenes Paar (Exziton).
  \item Zwei Elektronen mit entgegengesetzem Spin bilden
    ein Cooper-Paar (Supraleitung).
\end{itemize}

In allen drei Fällen entsteht ein gebundenes Paar aus einem
einzelnen Anregungszustand --- und in allen drei Fällen sind
Erhaltungsgrößen exakt erfüllt.

\begin{important}[Methodik dieses Dokuments]
Dieses Dokument trennt strikt zwei Ebenen:

\textbf{[STANDARD]}: Etablierte Physik, experimentell gesichert,
mit Primärquellen belegt.
Diese Aussagen sind unabhängig von FFGFT gültig.

\textbf{[FFGFT]}: Interpretation im FFGFT-Rahmen.
Diese Aussagen sind FFGFT-spezifisch und
sollten als solche gelesen werden.

Ziel ist es, die strukturelle Analogie zwischen den drei Systemen
rigoros zu belegen --- und dann präzise anzugeben,
wo die FFGFT-Deutung über die Standardphysik hinausgeht.
\end{important}

% ============================================================
\section{System 1: Parametrische Photonenaufspaltung (SPDC)}
\label{sec:spdc}
% ============================================================

\subsection{[STANDARD] Etablierte Physik der SPDC}

\textbf{Spontaneous Parametric Down-Conversion} (SPDC) ist ein
nichtlinearer optischer Prozess: ein Photon mit Frequenz $\omega_0$
wird in einem nichtlinearen Kristall in zwei Photonen
niedrigerer Frequenz umgewandelt~\cite{Burnham1970}:
\begin{equation}
  \omega_0 \;\longrightarrow\; \omega_1 + \omega_2,
  \quad \text{mit } \omega_1 + \omega_2 = \omega_0.
\end{equation}

Es gelten exakt zwei Erhaltungsgesetze:

\textbf{Energieerhaltung:}
\begin{equation}
  \hbar\omega_0 \;=\; \hbar\omega_1 + \hbar\omega_2.
\end{equation}

\textbf{Impulserhaltung (Phasenanpassung):}
\begin{equation}
  \mathbf{k}_0 \;=\; \mathbf{k}_1 + \mathbf{k}_2.
\end{equation}

Für Typ-II-SPDC haben die entstehenden Photonen
\textbf{entgegengesetzte Polarisation} und sind \textbf{verschränkt}.
Der Gesamtdrehimpuls des Systems ist erhalten~\cite{Kwiat1995}.

\begin{keyresult}[STANDARD: Was SPDC direkt zeigt]
\begin{enumerate}
  \item Ein höherenergetischer Zustand teilt sich in zwei
    niedrigenergetische auf.
  \item Energie und Impuls sind exakt erhalten.
  \item Die entgegengesetzte Polarisation der Photonen
    ist keine Zusatzannahme --- sie folgt aus der
    Drehimpulserhaltung.
  \item Die entstehenden Photonen sind \textbf{echte Photonen}
    --- keine Quasi-Teilchen, keine Rechenhilfen.
    \emph{[FFGFT-Vorbehalt: siehe unten --- Photonen sind in FFGFT
    keine Punkt-Teilchen sondern ausgedehnte Wellenpakete.]}
\end{enumerate}
\end{keyresult}

\subsection{[FFGFT] Interpretation der SPDC}

\begin{foundation}[FFGFT: SPDC als Torsionsaufspaltung --- Photonen als Wellenpakete]
In FFGFT trägt jedes Photon einen Drehimpuls als
topologische Eigenschaft der elektromagnetischen Feldtorsion
(Dokument~174, Abschnitt~3).

\textbf{Wichtige FFGFT-Grundposition: Photonen sind Wellenpakete.}

In der Standardtheorie (QED) ist das Photon ein quantisiertes
Anregungsquant des elektromagnetischen Feldes --- mathematisch
ein Fock-Zustand $\ket{1_\mathbf{k}}$.
Es wird oft als Punkt-Teilchen behandelt.

In FFGFT ist das Photon \textbf{kein Punkt-Teilchen},
sondern ein \textbf{ausgedehntes Wellenpaket}:
eine räumlich begrenzte, kohärente Torsionswelle des
elektromagnetischen Feldes mit definierter Frequenz $\omega$,
Wellenvektors $\mathbf{k}$ und Polarisation.

\textbf{Warum das für SPDC wichtig ist:}
Die Aufspaltung $\omega_0 \to \omega_1 + \omega_2$ ist
in FFGFT keine Zerlegung eines Punkt-Teilchens,
sondern die \textbf{Aufspaltung eines Wellenpakets}
in zwei Wellenpakete niedrigerer Frequenz im nichtlinearen Kristall.
Das ist physikalisch direkter --- das Licht \emph{ist} eine Welle,
und die Aufspaltung ist ein Wellenprozess.

\begin{center}
\begin{tabular}{p{5cm}p{6cm}}
\toprule
\textbf{Standardtheorie (QED)} & \textbf{FFGFT} \\
\midrule
Photon = Fock-Zustand $\ket{1_\mathbf{k}}$ &
  Photon = ausgedehntes Torsionswellenpaket \\
Punkt-Teilchen-Bild für viele Berechnungen &
  Kein Punkt-Teilchen; immer Wellenpaket \\
Aufspaltung = Vernichtung + Erzeugung von Quanten &
  Aufspaltung = Wellenprozess im nichtlinearen Medium \\
Verschränkung = Quantenkorrelation der Fock-Zustände &
  Verschränkung = topologische Torsionskorrelation \\
\bottomrule
\end{tabular}
\end{center}

\textbf{Konsequenz für die Beschreibung ``echtes Photon'':}
``Echt'' bedeutet in FFGFT: keine Quasi-Teilchen-Näherung,
keine Vielteilchen-Umformulierung --- aber kein Punktteilchen.
Ein SPDC-Photon ist ein echtes ausgedehntes Wellenpaket
des Torsionsfeldes mit klar definierten Eigenschaften
($\omega$, $\mathbf{k}$, Polarisation, räumliche Ausdehnung).

\textbf{Erhaltungsgröße in FFGFT:}
Die topologische Torsionsladung ist additiv und erhalten.
Das ist keine neue Behauptung --- es ist die FFGFT-Formulierung
der bekannten Drehimpulserhaltung, jetzt für Wellenpakete
statt für Punkt-Teilchen formuliert.
\end{foundation}

% ============================================================
\section{System 2: Das Exziton --- Elektron und Loch}
\label{sec:exziton}
% ============================================================

\subsection{[STANDARD] Was ein Loch wirklich ist}

Hier liegt der konzeptuelle Kern.

\begin{important}[STANDARD: Das Loch ist kein echtes Teilchen]
Ein \textbf{Loch} im Valenzband eines Halbleiters ist
\textbf{kein eigenständiges fundamentales Teilchen} ---
es ist ein \textbf{Quasi-Teilchen}, das die kollektive
Reaktion des Vielteilchensystems auf ein fehlendes Elektron
beschreibt~\cite{Kittel2004,AshcroftMermin1976}.

Präzise formuliert (Standardlehrbuch-Definition):
\begin{itemize}
  \item Das Valenzband ist im Grundzustand vollständig besetzt.
  \item Wenn ein Elektron fehlt (durch Absorption eines Photons),
    verhält sich das Vielteilchensystem \emph{so als wäre}
    ein positiv geladenes Teilchen mit positiver effektiver Masse
    vorhanden.
  \item Dieses fiktive Teilchen heißt ``Loch''.
  \item Es hat: positive Ladung $+e$, effektive Masse $m_h^*$,
    und einen Spin der dem des fehlenden Elektrons entgegengesetzt ist.
\end{itemize}

Das Loch ist mathematisch äquivalent zur Beschreibung
in Elektronen-Sprache --- aber es ist ein Konstrukt der
Vielteilchennäherung, kein echtes Antiteilchen
(kein Positron!).
\end{important}

Das echte Antiteilchen des Elektrons ist das \textbf{Positron} ---
das entsteht bei der Paarproduktion ($\gamma \to e^- + e^+$)
aus echten Photonen hoher Energie. Das Loch im Halbleiter
ist strukturell verschieden: es existiert nur im Kontext
des Valenzbandes des Materials.

\subsection{[STANDARD] Das Exziton}

Das \textbf{Exziton} ist ein gebundener Zustand aus einem
Leitungselektron und einem Valenzband-Loch, gebunden durch
Coulomb-Wechselwirkung~\cite{Frenkel1931,Wannier1937}:

\textbf{Frenkel-Exziton} (lokalisiert, Bindungsenergie 0,1--1 eV):
Das Elektron-Loch-Paar ist auf ein Molekül oder wenige
Gitterplätze beschränkt. Typisch in Isolatoren und
organischen Halbleitern~\cite{Frenkel1931}.

\textbf{Wannier-Mott-Exziton} (delokalisiert, Bindungsenergie $\sim$meV--100\,meV):
Das Elektron-Loch-Paar erstreckt sich über viele Gitterplätze ---
wasserstoffähnliche Beschreibung mit effektiven Massen und
dielektrischer Abschirmung~\cite{Wannier1937}. Für CdSe gilt:
\begin{equation}
  a_X \;=\; \frac{\varepsilon_r \hbar^2}{\mu e^2} \;\approx\; 5\text{--}6\,\text{nm},
\end{equation}
wobei $\mu$ die reduzierte Masse des Elektron-Loch-Paares ist.

\begin{keyresult}[STANDARD: Was beim Exziton erhalten ist]
\begin{enumerate}
  \item \textbf{Energie:} Die Anregungsenergie des Photons
    entspricht der Summe aus Bandlücke minus Exziton-Bindungsenergie.
  \item \textbf{Impuls:} Der Gesamtimpuls des Exzitons ist erhalten.
  \item \textbf{Spin:} Elektronen-Spin und Loch-Spin addieren sich
    zum Gesamtspin des Exzitons (Singulett $S=0$ oder Triplett $S=1$).
  \item \textbf{Ladung:} Das Exziton ist elektrisch neutral
    (Elektron $-e$ + Loch $+e = 0$).
\end{enumerate}
\end{keyresult}

\begin{important}[STANDARD: Kritischer Unterschied zur SPDC]
Beim Exziton entstehen \textbf{kein echtes Teilchen und kein echtes
Antiteilchen} --- sondern ein Elektron im Leitungsband und
ein Quasi-Teilchen (Loch) im Valenzband.

Die Beschreibung ``Elektron + Loch'' ist mathematisch
äquivalent zu einer Beschreibung in reinen Elektronen ---
das Loch ist eine \textbf{Rechenhilfe der Vielteilchentheorie},
die das komplizierte Vielteilchenproblem vereinfacht.

Auf dieser Ebene ist es korrekt zu sagen:
Ein Exziton besteht aus einem fehlenden Elektron im Valenzband
und einem überschüssigen Elektron im Leitungsband ---
also aus \textbf{einer Änderung in der Elektronenverteilung},
nicht aus zwei verschiedenen Teilchensorten.
\end{important}

\subsection{[FFGFT] Interpretation des Exzitons}

\begin{foundation}[FFGFT: Exziton als Torsion + Anti-Torsion]
In FFGFT (Dokument~173) ist jedes Elektron eine
stabile Torsionskonfiguration des Grundfeldes.
Das Loch --- die Abwesenheit eines Elektrons --- ist
in FFGFT-Sprache eine \textbf{Anti-Torsion}:
die geometrisch komplementäre Konfiguration zur Elektronen-Torsion.

Diese Beschreibung ist konsistent mit der Standardphysik:
Das Loch hat die entgegengesetzte Ladung und den
entgegengesetzten Spin des fehlenden Elektrons ---
genau wie eine Anti-Torsion die entgegengesetzte
Torsionswindung zur Elektronen-Torsion hätte.

\textbf{Das Exziton in FFGFT:}
Ein Torsions-Anti-Torsions-Paar --- gebunden durch
die geometrische Komplementarität der Feldkonfigurationen.
Das ist die FFGFT-Formulierung der Coulomb-Bindung im Exziton.

\textbf{Wichtige Einschränkung:}
Diese Beschreibung ist eine FFGFT-Interpretation,
nicht eine von der Standardphysik unabhängige Aussage.
Die Standardphysik beschreibt dasselbe Phänomen vollständig
und korrekt in Elektronen-Sprache.
\end{foundation}

% ============================================================
\section{System 3: Das Cooper-Paar --- zwei Elektronen}
\label{sec:cooper}
% ============================================================

\subsection{[STANDARD] Was ein Cooper-Paar ist}

Das \textbf{Cooper-Paar} ist fundamentell verschieden
von Exziton und SPDC. Das ist entscheidend:

\begin{keyresult}[STANDARD: Cooper-Paar = zwei echte Elektronen]
Ein Cooper-Paar besteht aus \textbf{zwei echten Elektronen}
mit entgegengesetzem Spin und entgegengesetzem Impuls:
\begin{equation}
  (\mathbf{k}, \uparrow) + (-\mathbf{k}, \downarrow)
  \;\longrightarrow\; \text{Cooper-Paar}.
\end{equation}
\textit{Kein} Loch, \textit{kein} Quasi-Teilchen, \textit{keine} fehlende Ladung.
Zwei echte Elektronen~\cite{Cooper1956,BCS1957}.
\end{keyresult}

\textbf{Wie entsteht die Anziehung?} (BCS-Mechanismus)

Elektronen stoßen sich normalerweise ab (Coulomb).
Die Paarung entsteht durch einen \textbf{indirekten Mechanismus}
über Gitterschwingungen (Phononen)~\cite{BCS1957}:

\begin{enumerate}
  \item Elektron~1 ($\mathbf{k},\uparrow$) bewegt sich durch
    das Kristallgitter und zieht positive Ionen leicht zu sich.
  \item Das deformierte Gitter erzeugt eine lokale
    positive Ladungserhöhung --- die nach dem Elektron
    ``nachschwingt'' (verzögerte Wechselwirkung).
  \item Elektron~2 ($-\mathbf{k},\downarrow$) wird von
    dieser Ladungserhöhung angezogen.
  \item Der Nettoeeffekt: eine schwache, indirekte Anziehung
    zwischen den zwei Elektronen --- vermittelt durch Phononen.
\end{enumerate}

Diese Anziehung ist \textbf{sehr schwach} (Bindungsenergie
$\sim 2\Delta \approx 3{,}5\,k_BT_c$).
Die Phonon-vermittelte Anziehung existiert bei jeder Temperatur ---
aber bei hoher Temperatur übersteigt die thermische Energie
die Bindungsenergie ($k_BT > \Delta$) und das
\textbf{Cooper-Paar wird aufgebrochen}~\cite{BCS1957}.
Erst bei tiefen Temperaturen ($k_BT < \Delta$) ist
die Paarung stabil.

\begin{important}[STANDARD: Warum tiefe Temperatur notwendig ist]
Die Cooper-Paar-Bindungsenergie liegt typischerweise
bei $\Delta \sim 1$--$10$\,meV.
Bei $T = 300$\,K ist $k_BT \approx 26$\,meV.
Die thermische Energie übersteigt die Bindungsenergie
bei Raumtemperatur um den Faktor~3--30.
Das Cooper-Paar wird thermisch aufgebrochen.

Erst unterhalb der kritischen Temperatur $T_c$
(typisch: 1--40\,K für konventionelle Supraleiter,
bis 135\,K für Hochtemperatur-Supraleiter)
überwiegt die Paarungswechselwirkung.

Das ist konsistent mit Dokument~173 (Bedingung~2):
$\Delta E > k_BT$ muss erfüllt sein.
Das Cooper-Paar sitzt auf einer anderen $\xi$-Stufe
als der Permanentmagnet --- seine Bindungsenergie
ist viel kleiner, also seine Stabilitätsgrenze
viel tiefer auf der Temperaturskala.
\end{important}

\subsection{[STANDARD] Cooper-Paar als Boson}

Das Cooper-Paar hat Gesamtspin $S = 0$ (Singulett, s-Welle)
oder $S = 1$ (Triplett, unkonventionelle Supraleiter).
Damit ist es ein zusammengesetztes Boson ---
es kann den Bose-Einstein-Statistiken gehorchen
und in den Grundzustand kondensieren (BCS-Kondensat)~\cite{BCS1957}.

\begin{center}
\resizebox{\textwidth}{!}{%
\begin{tabular}{p{3cm}p{3cm}p{3cm}p{3.5cm}}
\toprule
\textbf{Eigenschaft} & \textbf{Einzelelektron} &
\textbf{Cooper-Paar} & \textbf{Kommentar} \\
\midrule
Teilchenart    & Fermion & zusammenges. Boson &
  Gesamtspin ganzzahlig \\
Ladung         & $-e$ & $-2e$ & Messbar (Flux-Quantisierung) \\
Spin           & $\pm\frac{1}{2}$ & $0$ oder $1$ &
  Entgegengesetzt gepaart \\
Impuls         & $\mathbf{k}$ & $\mathbf{k} + (-\mathbf{k}) = 0$ &
  Schwerpunktsimpuls Null \\
Bindungsenergie & --- & $\sim 1$--$10$\,meV &
  Sehr schwach, $T_c < 40$\,K (konv.) \\
\bottomrule
\end{tabular}
}
\end{center}

\subsection{[FFGFT] Interpretation des Cooper-Paares}

\begin{foundation}[FFGFT: Cooper-Paar als entgegengesetzte Torsionskonfigurationen]
In FFGFT haben die zwei Elektronen des Cooper-Paares
entgegengesetzte Torsionskonfigurationen:
$(\mathbf{k}, \uparrow)$ trägt Torsion in eine Richtung,
$(-\mathbf{k}, \downarrow)$ trägt Torsion in die entgegengesetzte.

Das Paar hat damit:
\begin{itemize}
  \item Gesamttorsion Null (Spin-Singulett)
  \item Gesamtimpuls Null ($\mathbf{k} + (-\mathbf{k}) = 0$)
\end{itemize}

Das ist ein \textbf{topologisch neutraler Zustand} ---
geometrisch stabiler als zwei freie, getrennte Torsionskonfigurationen.

\textbf{Warum tiefe Temperatur nötig ist (FFGFT-Deutung):}
Die Bindung entsteht nicht direkt aus der Torsionsgeometrie,
sondern vermittelt über Phononen (Gitterschwingungen).
Phononen sind kollektive Schwingungsmoden des Kristallgitters ---
sie liegen auf einer anderen $\xi$-Stufe als das Elektron-Spin-System.
Die Phonon-vermittelte Kopplung ist deshalb eine
Querschnitts-Kopplung zwischen $\xi$-Stufen
(Dokument~175, Abschnitt über $\xi$-Hierarchie),
deren Stärke $\propto \xipar^{\Delta N}$ mit $\Delta N = 1$ ist.
Das erklärt die Schwäche der Bindungsenergie
im Vergleich zur direkten Spin-Spin-Kopplung.

\textbf{Kritische Einschränkung:}
Diese Deutung ist konsistent aber nicht bewiesen.
Die Standard-BCS-Theorie erklärt denselben Sachverhalt
vollständig ohne FFGFT-Konzepte.
\end{foundation}

% ============================================================
\section{Vergleich der drei Systeme}
\label{sec:vergleich}
% ============================================================

\subsection{[STANDARD] Strukturelle Gemeinsamkeiten}

\begin{center}
\resizebox{\textwidth}{!}{%
\begin{tabular}{p{2.5cm}p{3cm}p{3cm}p{3cm}p{2.5cm}}
\toprule
\textbf{System} & \textbf{Konstituenten} & \textbf{Bindungs\-energie} &
\textbf{Temperatur} & \textbf{Erhalten} \\
\midrule
\textbf{SPDC} &
  Photon~1 + Photon~2 &
  --- (keine Bindung) &
  Raumtemperatur &
  $E$, $\mathbf{p}$, $J_z$ \\[6pt]
\textbf{Exziton} &
  Elektron (Leitungsband) + Loch (Valenzband, Quasi-Teilchen) &
  $10$--$100$\,meV (Wannier) bis 1\,eV (Frenkel) &
  4\,K--300\,K je nach Typ &
  $E$, $\mathbf{p}$, $S$, Ladung \\[6pt]
\textbf{Cooper-Paar} &
  Elektron~1 ($\mathbf{k},\uparrow$) + Elektron~2 ($-\mathbf{k},\downarrow$) &
  $1$--$10$\,meV (konventionell) &
  $< T_c$ (1--40\,K) &
  $E$, $\mathbf{p}=0$, $S=0$ \\
\bottomrule
\end{tabular}
}
\end{center}

\textbf{Strukturelle Gemeinsamkeit [STANDARD]:}
In allen drei Fällen entstehen aus einem angeregten Zustand
zwei Objekte mit entgegengesetzten Quantenzahlen
(Spin, Impuls, Polarisation), sodass Erhaltungsgrößen exakt erfüllt sind.

\textbf{Kritischer Unterschied [STANDARD]:}
\begin{itemize}
  \item \textbf{SPDC:} zwei echte Photonen, keine Bindungsenergie,
    kein Temperaturlimit.
  \item \textbf{Exziton:} echtes Elektron + Quasi-Teilchen (Loch),
    Coulomb-Bindung, temperatursensitiv (Wannier-Mott: $T \ll T_\text{Bind}/k_B$).
  \item \textbf{Cooper-Paar:} zwei echte Elektronen, Phonon-vermittelte
    Bindung, sehr temperatursensitiv ($T < T_c$).
\end{itemize}

Das Loch beim Exziton ist \textbf{kein} Analogon zum zweiten Photon
bei SPDC --- es ist ein Quasi-Teilchen ohne eigenständige
fundamentale Existenz außerhalb des Vielteilchensystems.

\subsection{[STANDARD] Die Frage des Benutzers: ``Loch = zwei Elektronen?''}

Das ist eine präzise und wichtige Frage.

Die Antwort ist: \textbf{technisch ja, konzeptuell nein}.

\textbf{Technisch ja:}
Die Exziton-Beschreibung ``Elektron + Loch'' ist mathematisch
äquivalent zu einer Beschreibung ausschließlich in Elektronen.
In der Vielteilchentheorie ist das Loch definiert als:
\begin{equation}
  \ket{\text{Loch}} \;\equiv\; \ket{\text{Fermi-See ohne 1 Elektron bei }
  \mathbf{k}, \downarrow}.
\end{equation}
Das Loch \emph{ist} das fehlende Elektron in Verkleidung.

\textbf{Konzeptuell nein:}
Die Loch-Beschreibung ist nicht bloß eine Umbenennung ---
sie ist eine effiziente Reorganisation der Vielteilchen-Zustände,
die viele Phänomene (Leitfähigkeit, Magnetotransport,
optische Übergänge) viel einfacher beschreibt als die
reine Elektronen-Sprache.
Das Loch \emph{verhält sich} wie ein Teilchen --- es hat
eine effektive Masse, einen Spin, einen Impuls.
Es ist ein nützliches Konstrukt mit realem physikalischen Inhalt,
auch wenn es kein fundamentales Teilchen ist.

\begin{keypoint}[Loch: nützliches Konstrukt, kein fundamentales Teilchen]
Das Loch im Halbleiter ist nicht:
\begin{itemize}
  \item Ein echtes Antiteilchen (kein Positron)
  \item Eine neue Fundamentalstruktur der Natur
  \item Unabhängig von der Elektronen-Beschreibung
\end{itemize}

Das Loch ist:
\begin{itemize}
  \item Ein Quasi-Teilchen: eine emergente Anregung des
    Vielteilchensystems
  \item Mathematisch: die Absence eines Elektrons im
    Valenzband, umgeschrieben als positiv geladenes Teilchen
  \item Physikalisch sinnvoll: es beschreibt reale,
    messbare Phänomene korrekt
\end{itemize}

Das \textbf{Cooper-Paar} besteht hingegen aus
\textbf{zwei echten Elektronen} --- kein Loch,
kein Quasi-Teilchen. Das ist der fundamentale Unterschied.
\end{keypoint}

% ============================================================
\section{[FFGFT] Das gemeinsame Prinzip: Torsionserhaltung}
\label{sec:torsionserhaltung}
% ============================================================

\begin{foundation}[FFGFT: Erhaltung der Gesamttorsion als einheitliches Prinzip]
In FFGFT ist das verbindende Prinzip hinter allen drei Systemen
die \textbf{Erhaltung der Gesamttorsion}:

\begin{center}
\resizebox{\textwidth}{!}{%
\begin{tabular}{p{2.5cm}p{4cm}p{4cm}p{3cm}}
\toprule
\textbf{System} & \textbf{Vor dem Prozess} &
\textbf{Nach dem Prozess} & \textbf{Erhaltenes} \\
\midrule
SPDC &
  Torsion $\omega_0$, Drehimpuls $J_z$ &
  Torsion $\omega_1$ (+) + Anti-Torsion $\omega_2$ (--) &
  Torsionsladung: $J_z = J_1 + J_2$ \\[6pt]
Exziton &
  Valenzband-Torsion (volles Band) &
  Leitungsband-Torsion (Elektron) + Valenzband-Anti-Torsion (Loch) &
  Torsionsladung: $S_\text{total}$, $\mathbf{p}_\text{total}$ \\[6pt]
Cooper-Paar &
  Fermi-See-Torsion & Paar-Torsion ($\mathbf{k},\uparrow$) + ($-\mathbf{k},\downarrow$) &
  Torsionsladung: $S = 0$, $\mathbf{p} = 0$ \\
\bottomrule
\end{tabular}
}
\end{center}

\textbf{Die gemeinsame FFGFT-Sprache:}
In allen drei Fällen entsteht aus einer Torsionskonfiguration
ein Paar aus Torsion und Anti-Torsion, sodass die
Gesamttorsionsladung erhalten bleibt.

Das ist in FFGFT kein Zufall --- es ist das Spiegelbild
der bekannten Erhaltungssätze (Energie, Impuls, Drehimpuls)
in der Sprache des Torsionsfeldes.

\textbf{Was die FFGFT nicht behauptet:}
FFGFT behauptet \emph{nicht}, dass Exziton-Loch und
Cooper-Elektron und SPDC-Photon dieselben Objekte sind ---
sie sind physikalisch verschieden.
FFGFT behauptet, dass ihre \emph{strukturelle Ähnlichkeit}
einen gemeinsamen geometrischen Ursprung hat:
die Erhaltung der topologischen Torsionsladung.
\end{foundation}

% ============================================================
\section{Konsistenz mit dem Temperaturargument aus Dokument~173}
\label{sec:temperatur}
% ============================================================

Die drei Systeme unterscheiden sich dramatisch in ihrer
Temperaturstabilität. Das ist [STANDARD] und direkt erklärbar:

\begin{center}
\resizebox{\textwidth}{!}{%
\begin{tabular}{p{2.8cm}p{2.5cm}p{3cm}p{4cm}}
\toprule
\textbf{System} & \textbf{Bindungsenergie} &
\textbf{Stabilitätstemperatur} & \textbf{Ursache} \\
\midrule
Photon (SPDC) &
  keine (freie Photonen) &
  Raumtemperatur &
  Keine thermische Dissipation \\[6pt]
Frenkel-Exziton &
  0{,}1--1\,eV &
  Raumtemperatur möglich &
  $\Delta E \gg k_BT(300\,\text{K}) = 26$\,meV \\[6pt]
Wannier-Exziton (CdSe) &
  $\sim 15$--20\,meV &
  $< 50$\,K &
  $\Delta E \sim k_BT$ bei Raumtemperatur \\[6pt]
Cooper-Paar (konvent.) &
  1--10\,meV &
  $T < T_c = 1$--40\,K &
  $\Delta E \ll k_BT(300\,\text{K})$ \\
\bottomrule
\end{tabular}
}
\end{center}

\begin{keypoint}[Konsistenz mit Bedingung 2 aus Dokument~173]
Die Temperaturgrenzen folgen direkt aus Bedingung~2
(Dokument~173): $\Delta E > k_BT$.
Je schwächer die Bindungsenergie, desto tiefer muss
die Temperatur sein --- das ist kein FFGFT-Argument,
das ist Standardphysik.

Das Cooper-Paar braucht tiefe Temperatur nicht wegen
seiner Elektronennatur, sondern wegen seiner
\textbf{sehr schwachen Bindungsenergie},
die aus dem indirekten Phonon-Mechanismus folgt.
\end{keypoint}

% ============================================================
\section{Zusammenfassung}
\label{sec:zusammenfassung-179}
% ============================================================

\begin{conclusionBox}[Fazit: Was sicher ist und was FFGFT-Deutung ist]

\textbf{[STANDARD] Gesicherte Aussagen:}

\textbf{1. Das Loch ist ein Quasi-Teilchen, kein echtes Antiteilchen.}\\
Es ist die mathematische Umformulierung eines fehlenden Elektrons
im Valenzband. ``Exziton besteht aus zwei Elektronen'' ist
technisch korrekt --- aber die Loch-Beschreibung ist effizienter
und physikalisch sinnvoll.

\textbf{2. Das Cooper-Paar besteht aus zwei echten Elektronen.}\\
Entgegengesetzter Spin und Impuls. Phonon-vermittelte Bindung.
Tiefe Temperatur ist nötig wegen der schwachen Bindungsenergie
($\Delta E \ll k_BT_\text{Raumtemperatur}$).

\textbf{3. SPDC erzeugt zwei echte Photonen --- in FFGFT: zwei Wellenpakete.}\\
Energie, Impuls und Drehimpuls sind exakt erhalten.
Keine Bindungsenergie --- die Photonen sind frei.
In FFGFT sind Photonen keine Punkt-Teilchen sondern
ausgedehnte Torsionswellenpakete --- die Aufspaltung
ist ein Wellenprozess, kein Zerfall eines Punktteilchens.

\textbf{4. Die strukturelle Ähnlichkeit ist real:}\\
In allen drei Fällen entstehen zwei Objekte mit
entgegengesetzten Quantenzahlen aus einem Ausgangszustand,
sodass Erhaltungsgrößen exakt erfüllt sind.

\medskip

\textbf{[FFGFT] Interpretation:}

\textbf{5. Torsionserhaltung als einheitliches Prinzip.}\\
Die bekannten Erhaltungssätze (Energie, Impuls, Drehimpuls)
entsprechen in FFGFT der Erhaltung der topologischen
Torsionsladung. Die strukturelle Ähnlichkeit der drei Systeme
hat damit einen gemeinsamen geometrischen Ursprung ---
ohne dass die drei Systeme dasselbe wären.

\textbf{6. Das Loch als Anti-Torsion.}\\
Konsistent mit der FFGFT-Beschreibung des Exzitons
aus Dokument~173 --- aber über die Standardphysik hinausgehend
und nicht unabhängig bewiesen.
\end{conclusionBox}

% ============================================================

\begin{thebibliography}{99}

% --- SPDC ---
\bibitem{Burnham1970}
Burnham, D.\,C.\ \& Weinberg, D.\,L.
\textit{Observation of simultaneity in parametric production of
optical photon pairs.}
Phys.\ Rev.\ Lett.\ \textbf{25}, 84--87 (1970).

\bibitem{Kwiat1995}
Kwiat, P.\,G.\ et al.
\textit{New high-intensity source of polarization-entangled photon pairs.}
Phys.\ Rev.\ Lett.\ \textbf{75}, 4337--4341 (1995).

% --- Exziton ---
\bibitem{Frenkel1931}
Frenkel, J.
\textit{On the transformation of light into heat in solids. I \& II.}
Phys.\ Rev.\ \textbf{37}, 17--44 und 1276--1294 (1931).

\bibitem{Wannier1937}
Wannier, G.\,H.
\textit{The structure of electronic excitation levels in
insulating crystals.}
Phys.\ Rev.\ \textbf{52}, 191--197 (1937).

% --- Halbleiter-Grundlagen (Loch als Quasi-Teilchen) ---
\bibitem{Kittel2004}
Kittel, C.
\textit{Introduction to Solid State Physics.}
8th ed.\ Wiley (2004). Kapitel~8: Semiconductor Crystals.

\bibitem{AshcroftMermin1976}
Ashcroft, N.\,W.\ \& Mermin, N.\,D.
\textit{Solid State Physics.}
Holt, Rinehart and Winston (1976). Kapitel~12: The Semiclassical
Model of Electron Dynamics; Kapitel~28: Homogeneous Semiconductors.

% --- Cooper-Paar und BCS ---
\bibitem{Cooper1956}
Cooper, L.\,N.
\textit{Bound electron pairs in a degenerate Fermi gas.}
Phys.\ Rev.\ \textbf{104}, 1189--1190 (1956).

\bibitem{BCS1957}
Bardeen, J., Cooper, L.\,N.\ \& Schrieffer, J.\,R.
\textit{Theory of superconductivity.}
Phys.\ Rev.\ \textbf{108}, 1175--1204 (1957).

% --- Exziton in Quantenpunkten ---
\bibitem{Woggon1997}
Woggon, U.
\textit{Optical Properties of Semiconductor Quantum Dots.}
Springer Tracts in Modern Physics, Vol.\ 136 (1997).

\end{thebibliography}
